\section{Parte A - Control de rumbo con asignación de polos}

\subsection{Modelo en espacio de estados}

Con el vector de estados $x = [v\ r\ \psi]^T$, la entrada $u_c = \tau_r$ y la salida $y = \psi$, el modelo lineal puede escribirse como

\begin{equation}
\dot{x} = A x + B u_c, \qquad y = Cx + Du_c,
\end{equation}

donde

\begin{equation}
A =
\begin{bmatrix}
-0.25 & -2 & 0 \\
0.08 & -0.35 & 0 \\
0 & 1 & 0
\end{bmatrix},
\quad
B =
\begin{bmatrix}
0 \\
0.62 \\
0
\end{bmatrix},
\quad
C =
\begin{bmatrix}
0 & 0 & 1
\end{bmatrix},
\quad
D = 0.
\end{equation}

\subsection{Controlabilidad del sistema original}

La matriz de controlabilidad del sistema original se define como

\begin{equation}
\mathcal{C} =
\begin{bmatrix}
B & AB & A^2B
\end{bmatrix}
=
\begin{bmatrix}
0 & -1.2400 & 0.7440 \\
0.6200 & -0.2170 & -0.0233 \\
0 & 0.6200 & -0.2170
\end{bmatrix}.
\label{eq:controlabilidad-original}
\end{equation}

Como $\operatorname{rango}(\mathcal{C})=3$, la matriz tiene rango completo e igual a la dimensión del sistema original. Por lo tanto, es posible diseñar un controlador por realimentación completa de estados para el modelo de rumbo.

\subsection{Sistema aumentado con acción integral}

Para eliminar el error estacionario frente a referencias constantes o por tramos se incorpora el estado integral

\begin{equation}
\dot{\xi} = \psi_{ref} - \psi.
\end{equation}

El vector aumentado es $x_a = [v\ r\ \psi\ \xi]^T$ y el sistema queda

\begin{equation}
\dot{x}_a = A_a x_a + B_a \tau_r + E_a \psi_{ref},
\end{equation}

con

\begin{equation}
A_a =
\begin{bmatrix}
-0.25 & -2 & 0 & 0 \\
0.08 & -0.35 & 0 & 0 \\
0 & 1 & 0 & 0 \\
0 & 0 & -1 & 0
\end{bmatrix},
\quad
B_a =
\begin{bmatrix}
0 \\
0.62 \\
0 \\
0
\end{bmatrix},
\quad
E_a =
\begin{bmatrix}
0 \\
0 \\
0 \\
1
\end{bmatrix}.
\end{equation}

\subsection{Controlabilidad del sistema aumentado}

La matriz de controlabilidad del sistema aumentado es

\begin{equation}
\mathcal{C}_a =
\begin{bmatrix}
B_a & A_aB_a & A_a^2B_a & A_a^3B_a
\end{bmatrix}
=
\begin{bmatrix}
0 & -1.2400 & 0.7440 & -0.1395 \\
0.6200 & -0.2170 & -0.0233 & 0.0677 \\
0 & 0.6200 & -0.2170 & -0.0233 \\
0 & 0 & -0.6200 & 0.2170
\end{bmatrix}.
\label{eq:controlabilidad-aumentado}
\end{equation}

Como $\operatorname{rango}(\mathcal{C}_a)=4$, la matriz tiene rango completo e igual a la dimensión del sistema aumentado. Por lo tanto, es posible asignar arbitrariamente los polos del sistema aumentado.

\subsection{Diseño del controlador}

Se utiliza la ley de control

\begin{equation}
\tau_r = -K_a x_a = -K_x x - k_i \xi,
\qquad
K_a = [k_v\ k_r\ k_\psi\ k_i].
\end{equation}

Primero se eligieron los polos deseados del sistema aumentado y luego se calculó la ganancia $K_a$. Para ello se utilizó la función \texttt{control.place}. Esta función es un wrapper de \texttt{scipy.signal.place\_poles}, que implementa el algoritmo de Tits-Yang para asignación robusta de polos por realimentación de estados.

Los polos asignados y la ganancia obtenida se muestran, en ese orden, en las Tablas~\ref{tab:parte-a-polos} y~\ref{tab:parte-a-ganancias}.

\begin{table}[H]
\centering
\caption{Polos asignados para el sistema aumentado de la Parte A.}
\label{tab:parte-a-polos}
\begin{tabular}{cccc}
\toprule
$p_1$ & $p_2$ & $p_3$ & $p_4$ \\
\midrule
-1.500 & -2.000 & -2.500 & -3.000 \\
\bottomrule
\end{tabular}
\end{table}

\begin{table}[H]
\centering
\caption{Ganancia del controlador por asignación de polos de la Parte A.}
\label{tab:parte-a-ganancias}
\begin{tabular}{rrrr}
\toprule
$k_v$ & $k_r$ & $k_\psi$ & $k_i$ \\
\midrule
-174.518 & 13.548 & -304.839 & -145.161 \\
\bottomrule
\end{tabular}
\end{table}


\subsection{Simulación de seguimiento}

La simulación se realiza durante $T_f = 60\,\mathrm{s}$ utilizando la referencia por tramos pedida en la consigna. En la Figura~\ref{fig:parte-a-seguimiento} se muestra el seguimiento de rumbo.

\begin{figure}[H]
\centering
\includegraphics[width=0.86\textwidth]{parteA/seguimiento.png}
\caption{Seguimiento de referencia de rumbo para el controlador por asignación de polos.}
\label{fig:parte-a-seguimiento}
\end{figure}

\subsection{Error, estados y señal de control}

La Figura~\ref{fig:parte-a-error-estados-control} resume el error de seguimiento, los estados del sistema y la señal de control. Las métricas principales se muestran en la Tabla~\ref{tab:parte-a-metricas}.

\begin{figure}[H]
\centering
\includegraphics[width=0.92\textwidth]{parteA/error_estados_control.png}
\caption{Error de seguimiento, estados y señal de control para la Parte A.}
\label{fig:parte-a-error-estados-control}
\end{figure}

\begin{table}[H]
\centering
\caption{Métricas principales de la simulación de la Parte A.}
\label{tab:parte-a-metricas}
\begin{tabular}{lr}
\toprule
Métrica & Valor \\
\midrule
$\max |e(t)|$ & 1.9810 \\
$e(T_f)$ & 0.0000 \\
IAE & 9.2710 \\
ISE & 9.8294 \\
$\max |\tau_r(t)|$ & 18.5989 \\
RMS de $\tau_r(t)$ & 1.9077 \\
$\max |v(t)|$ & 5.8903 \\
$\max |r(t)|$ & 4.0537 \\
$\max |\psi(t)|$ & 3.0282 \\
\bottomrule
\end{tabular}
\end{table}


\subsection{Discusión de desempeño}

\pendiente{Discutir si la embarcación sigue la referencia, si el error estacionario es nulo, si la velocidad de respuesta es adecuada, la suavidad de $\tau_r$, los valores máximos requeridos de $\tau_r$, el efecto esperado al mover los polos y las limitaciones físicas del actuador.}

\FloatBarrier
